% GENERATED FILE. Edit canonical lessons, then run npm run generate:books.
% canonical-source-hash: fnv1a64:7490e26ada9d2922
% canonical-lessons: SA-C358-acarah, SA-C358-silam, SA-C358-sajjanah, SA-C358-sanmargah, SA-C358-adriyate

\chapter{Good Conduct}
\label{ch:sa-good-conduct}
\hlchaptermodality{358}

\begin{chapteropening}
\textbf{By the end of this chapter:} \emph{I can talk about good conduct and good people.}

Complete the last lesson of chapter 358: I can talk about good conduct and good people.
\end{chapteropening}

\section[\texorpdfstring{\sk{आचारः} (ācāraḥ)}{आचारः (ācāraḥ)}]{\sk{आचारः} (ācāraḥ) — conduct; custom}

\label{lesson:SA-C358-acarah}



\begin{quote}
Before the new one: say the Sanskrit for a secretary, then the Sanskrit for directs.
\end{quote}

\subsection*{You'll want to know: \sk{आचारः}}
\textbf{\sk{आचारः}} — \emph{ācāraḥ} — \textquotedblleft{}conduct; custom\textquotedblright{}.

\begin{quote}
\sk{मीरायाः} \sk{आचारः} \sk{उत्तमः} \sk{अस्ति।}
\end{quote}

\begin{quote}
\emph{mīrāyāḥ ācāraḥ uttamaḥ asti.}
\end{quote}

\begin{quote}
Mira's conduct is excellent.
\end{quote}

A masculine noun.

\begin{cousinweb}
\textbf{\sk{आचारः}} (\emph{ācāraḥ}) is built on \emph{ā} and the root √car, \textquotedblleft{}to move\textquotedblright{}; MW gives \emph{ā-cāra} \textquotedblleft{}conduct, behaviour\textquotedblright{} and \textquotedblleft{}custom, practice\textquotedblright{} (MW).
\end{cousinweb}

\begin{grammarlens}[title={What you've built}]
One more word for giving reasons and opinions.
\end{grammarlens}

\subsection*{Your turn}
\begin{itemize}
\raggedright
  \item \emph{Say it:} \emph{ācāraḥ}
  \item \emph{Say it:} \emph{ācāraḥ}, once more
  \item \emph{Say it:} \emph{mīrāyāḥ ācāraḥ uttamaḥ asti.} — Mira's conduct is excellent.
  \item \emph{Recall:} say the Sanskrit for a secretary, then the Sanskrit for directs, then say \emph{ācāraḥ} again
\end{itemize}

\begin{tcolorbox}[breakable,colback=teal!4,colframe=teal!35!black,title={Before you move on}]
What does \sk{आचारः} mean? (\textquotedblleft{}Conduct; custom\textquotedblright{}.) Say it once more.
\end{tcolorbox}
\section[\texorpdfstring{\sk{शीलम्} (śīlam)}{शीलम् (śīlam)}]{\sk{शीलम्} (śīlam) — character; disposition}

\label{lesson:SA-C358-silam}



\begin{quote}
Before the new one: say the Sanskrit for directs, then the Sanskrit for conduct.
\end{quote}

\subsection*{You'll want to know: \sk{शीलम्}}
\textbf{\sk{शीलम्}} — \emph{śīlam} — \textquotedblleft{}character; disposition\textquotedblright{}.

\begin{quote}
\sk{बालस्य} \sk{शीलं} \sk{सुन्दरम्} \sk{अस्ति।}
\end{quote}

\begin{quote}
\emph{bālasya śīlaṃ sundaram asti.}
\end{quote}

\begin{quote}
The boy has a fine character.
\end{quote}

A neuter noun.

\begin{cousinweb}
\textbf{\sk{शीलम्}} (\emph{śīlam}) MW gives \emph{śīla} \textquotedblleft{}habit, custom, way of living, conduct, disposition\textquotedblright{} and gives no derivation, so its origin is uncertain (MW).
\end{cousinweb}

\begin{grammarlens}[title={What you've built}]
One more word for giving reasons and opinions.
\end{grammarlens}

\subsection*{Your turn}
\begin{itemize}
\raggedright
  \item \emph{Say it:} \emph{śīlam}
  \item \emph{Say it:} \emph{śīlam}, once more
  \item \emph{Say it:} \emph{bālasya śīlaṃ sundaram asti.} — The boy has a fine character.
  \item \emph{Recall:} say the Sanskrit for directs, then the Sanskrit for conduct, then say \emph{śīlam} again
\end{itemize}

\begin{tcolorbox}[breakable,colback=teal!4,colframe=teal!35!black,title={Before you move on}]
What does \sk{शीलम्} mean? (\textquotedblleft{}Character; disposition\textquotedblright{}.) Say it once more.
\end{tcolorbox}
\section[\texorpdfstring{\sk{सज्जनः} (sajjanaḥ)}{सज्जनः (sajjanaḥ)}]{\sk{सज्जनः} (sajjanaḥ) — a good person}

\label{lesson:SA-C358-sajjanah}



\begin{quote}
Before the new one: say the Sanskrit for conduct, then the Sanskrit for character.
\end{quote}

\subsection*{You'll want to know: \sk{सज्जनः}}
\textbf{\sk{सज्जनः}} — \emph{sajjanaḥ} — \textquotedblleft{}a good person\textquotedblright{}.

\begin{quote}
\sk{सज्जनः} \sk{सदा} \sk{सत्यं} \sk{वदति।}
\end{quote}

\begin{quote}
\emph{sajjanaḥ sadā satyaṃ vadati.}
\end{quote}

\begin{quote}
A good person tells the truth at all times.
\end{quote}

A masculine noun.

\begin{cousinweb}
\textbf{\sk{सज्जनः}} (\emph{sajjanaḥ}) is built from \emph{sat}, \textquotedblleft{}good\textquotedblright{}, and \emph{jana}, \textquotedblleft{}a person\textquotedblright{}; MW gives \emph{saj-jana} \textquotedblleft{}a good or virtuous or wise man\textquotedblright{} (MW).
\end{cousinweb}

\begin{grammarlens}[title={What you've built}]
One more word for giving reasons and opinions.
\end{grammarlens}

\subsection*{Your turn}
\begin{itemize}
\raggedright
  \item \emph{Say it:} \emph{sajjanaḥ}
  \item \emph{Say it:} \emph{sajjanaḥ}, once more
  \item \emph{Say it:} \emph{sajjanaḥ sadā satyaṃ vadati.} — A good person tells the truth at all times.
  \item \emph{Recall:} say the Sanskrit for conduct, then the Sanskrit for character, then say \emph{sajjanaḥ} again
\end{itemize}

\begin{tcolorbox}[breakable,colback=teal!4,colframe=teal!35!black,title={Before you move on}]
What does \sk{सज्जनः} mean? (\textquotedblleft{}A good person\textquotedblright{}.) Say it once more.
\end{tcolorbox}
\section[\texorpdfstring{\sk{सन्मार्गः} (sanmārgaḥ)}{सन्मार्गः (sanmārgaḥ)}]{\sk{सन्मार्गः} (sanmārgaḥ) — the right path}

\label{lesson:SA-C358-sanmargah}



\begin{quote}
Before the new one: say the Sanskrit for character, then the Sanskrit for a good person.
\end{quote}

\subsection*{You'll want to know: \sk{सन्मार्गः}}
\textbf{\sk{सन्मार्गः}} — \emph{sanmārgaḥ} — \textquotedblleft{}the right path\textquotedblright{}.

\begin{quote}
\sk{गुरुः} \sk{शिष्यं} \sk{सन्मार्गं} \sk{नयति।}
\end{quote}

\begin{quote}
\emph{guruḥ śiṣyaṃ sanmārgaṃ nayati.}
\end{quote}

\begin{quote}
The teacher leads his pupil on the right path.
\end{quote}

A masculine noun.

\begin{cousinweb}
\textbf{\sk{सन्मार्गः}} (\emph{sanmārgaḥ}) is built from \emph{sat}, \textquotedblleft{}good\textquotedblright{}, and \emph{mārga}, \textquotedblleft{}a path\textquotedblright{}; MW gives \emph{san-mārga} \textquotedblleft{}the right path\textquotedblright{} (MW).
\end{cousinweb}

\begin{grammarlens}[title={What you've built}]
One more word for giving reasons and opinions.
\end{grammarlens}

\subsection*{Your turn}
\begin{itemize}
\raggedright
  \item \emph{Say it:} \emph{sanmārgaḥ}
  \item \emph{Say it:} \emph{sanmārgaḥ}, once more
  \item \emph{Say it:} \emph{guruḥ śiṣyaṃ sanmārgaṃ nayati.} — The teacher leads his pupil on the right path.
  \item \emph{Recall:} say the Sanskrit for character, then the Sanskrit for a good person, then say \emph{sanmārgaḥ} again
\end{itemize}

\begin{tcolorbox}[breakable,colback=teal!4,colframe=teal!35!black,title={Before you move on}]
What does \sk{सन्मार्गः} mean? (\textquotedblleft{}The right path\textquotedblright{}.) Say it once more.
\end{tcolorbox}
\section[\texorpdfstring{\sk{आद्रियते} (ādriyate)}{आद्रियते (ādriyate)}]{\sk{आद्रियते} (ādriyate) — holds in regard; respects}

\label{lesson:SA-C358-adriyate}



\begin{quote}
Before the new one: say the Sanskrit for a good person, then the Sanskrit for the right path.
\end{quote}

\subsection*{You'll want to know: \sk{आद्रियते}}
\textbf{\sk{आद्रियते}} — \emph{ādriyate} — \textquotedblleft{}holds in regard; respects\textquotedblright{}.

\begin{quote}
\sk{छात्राः} \sk{गुरुम्} \sk{आद्रियन्ते।}
\end{quote}

\begin{quote}
\emph{chātrāḥ gurum ādriyante.}
\end{quote}

\begin{quote}
The students hold their teacher in regard.
\end{quote}

Past (he …): \sk{सः} \sk{गुरुम्} \sk{आदृतवान्}, \emph{saḥ gurum ādṛtavān}.

\begin{cousinweb}
\textbf{\sk{आद्रियते}} (\emph{ādriyate}) is built on \emph{ā} and the root √dṛ, \textquotedblleft{}to heed\textquotedblright{}; MW gives \emph{ā-√dṛ} \textquotedblleft{}to regard with attention\textquotedblright{} and \textquotedblleft{}to respect, to honour\textquotedblright{} (MW).
\end{cousinweb}

\begin{grammarlens}[title={What you've built}]
One more word for giving reasons and opinions.
\end{grammarlens}

\subsection*{Your turn}
\begin{itemize}
\raggedright
  \item \emph{Say it:} \emph{ādriyate}
  \item \emph{Say it:} \emph{ādriyate}, once more
  \item \emph{Say it:} \emph{chātrāḥ gurum ādriyante.} — The students hold their teacher in regard.
  \item \emph{Recall:} say the Sanskrit for a good person, then the Sanskrit for the right path, then say \emph{ādriyate} again
\end{itemize}

\begin{tcolorbox}[breakable,colback=teal!4,colframe=teal!35!black,title={Before you move on}]
What does \sk{आद्रियते} mean? (\textquotedblleft{}Holds in regard; respects\textquotedblright{}.) Say it once more.
\end{tcolorbox}
